\documentclass[12pt]{article}
\usepackage[utf8]{inputenc}
\usepackage[T1]{fontenc}
\usepackage{hyperref}
\title{Ura Language Resources for K-12 Curriculum Development}
\author{Compiled by Claude Sonnet 5 (claude-sonnet-5) for Andr\'as Kornai}
\date{}
\begin{document}
\maketitle

\section{Introduction}
Ura (ISO 639-3: uur) is an Austronesian language of Erromango, Vanuatu. It was thought extinct until Terry Crowley found a handful of elderly speakers in the 1990s: six fluent and a couple of dozen semi-fluent. One source considers it unlikely to be in use today, and I found no more recent count. UNESCO lists it as Critically Endangered. Documentation is purely academic; no teaching materials were found.

\section{Primary Educational Resources}
None found.

\section{Linguistic and Reference Materials}
\begin{itemize}
\item Terry Crowley, \emph{Ura: A disappearing language of southern Vanuatu}, Pacific Linguistics C-156 (1999), the main description. The book itself was not located online; a review is hosted at Project MUSE. \url{https://muse.jhu.edu/article/25417}
\item Terry Crowley, \emph{Ura}, LINCOM Languages of the World/Materials 240 (1998). It is listed in the Wikipedia article. \url{https://en.wikipedia.org/wiki/Ura_language_%28Vanuatu%29}
\item John Lynch's 1982 and 1983 papers on the Ura language of Erromango, listed in the Erromango language bibliography. \url{https://nthieberger.net/vanlangs/ERRO.HTM}
\end{itemize}

\section{Religious and Community Resources}
None found.

\section{Digital and Audio Resources}
The Endangered Languages Project has a profile of the language. No audio or corpus was identified. \url{https://endangeredlanguages.com/lang/2446}

\section{Curriculum Development Context}
A revitalization curriculum would rest mainly on Crowley's grammar and any recordings he made, whose location is unknown to me. The closely related Sye (Erromanga) language is not endangered and might serve as a comparison. Community consultation in Erromango would be essential.
\end{document}